\section{Assembling the actual segment spaces}
\label{sec:segment-assembly}

\status{This closes the actual segment-product injection and degree matching retained in SD2. With the shared integral path identification CG1 and the numerical bound RS3, it proves \texttt{RestLam} under the source's \texttt{Dat.Hyp}, with $C_{\rm La},C_k\ge1088$. Counting sufficiently many segment paths is a separate argument.}

\paragraph{LA1--LA2. The frame geometry and its scales.}
Write $k=r+K$, $f_{j,i}=\mathrm{fr}(j,i)$, $z_j=f_{j,r-1}$ and $s_j=f_{j-1,0}$.
Frame validity says that the $r$ selected frame letters are ordered, consecutive frames are separated, there are no selected letters between adjacent entries of a frame, and
$f_{j+J,i}=f_{j,i}+n'$.
Consequently segment $j$ has precisely the old-frame letters $f_{j-1,i}$ and the late selected letters
\[
 L_j=TT_{\mathbb Z}\cap(z_{j-1},z_j].
\]
For $0\le j<J$, these late sets partition the selected residue positions: their union lies in the half-open interval $(z_{-1},z_{-1}+n']$.
Old-frame positions likewise occur at most once per period. An old/late overlap is exactly a current-frame letter, shared with the next segment.

The apparently complicated coarse powers have a simple exact form. Put
\[
 H=\log_p\left\lfloor\frac{a}{8Yk^2N_s}\right\rfloor,
 \qquad A_j=\max\{z_j,H+s_j\}.
\]
Here $\log_p$ denotes the integer logarithm, with $\log_p0=0$.
The identity $\log_p(N/p^d)=\max(\log_pN-d,0)$, with integer division, gives
\[
 R_j=\max\{H-(z_j-s_j),0\},\qquad
 P_j(x)=p^{A_j-x}\ (x\le z_j),\qquad Q_j=p^{A_{j+1}-A_j}. \tag{LA.1}
\]
Both $z_j$ and $s_j$ increase strictly, so these exponents have the stated signs.
Also $A_{j+J}=A_j+n'$. Thus
\[
 P_{j+1}(x)=P_j(x)Q_j\quad(x\le z_j),
 \qquad \prod_{j=0}^{J-1}Q_j=p^{n'}. \tag{LA.2}
\]
These facts prove every condition of the source's \texttt{LamGeo}, including the overlap at the period boundary.

\paragraph{LA3. Saturation for the actual restricted spaces.}
Let $O$ be the source DVR with residue field $\mathbb f$.
Its segment lattice is exactly
\[
 M_j(b)=XX'_j(m_0,\theta_b,o_c)
 \cap\{\text{row degrees }b\}
 \cap\{m(f_{j,i},c)<Q_j\text{ for }c<r\},
 \qquad m_0=kN_s+1.
\]
CG1 identifies the first module with the span of a subset of the same integral paired path basis over $O$ and $\mathbb f$; hence it is saturated.
Its hypotheses hold because $m_0\ge0$ and
$\#\mathrm{letters}_j+r\le N_s+r\le K$.
The two remaining restrictions are coefficient-support conditions, so their intersections preserve saturation.
Fixed row degrees and finite letter support make $M_j(b)$ finite free.
An $O$-basis stays independent after reduction: a lifted relation
$\sum c_i b_i=\varpi P$ puts $P$ back in $M_j(b)$ by saturation, and basis coordinates imply $\varpi\mid c_i$ for every $i$.
This establishes the actual \texttt{RestLamA} input, rather than assuming that reduction preserves an arbitrary lattice basis.

\paragraph{LA4. Recovering every label and segment monomial.}
An allowed label $\ell$ satisfies, at the unique late representative $x\in L_j$ of $t$,
\[
 0\le\ell_{t,c}<P_j(x),\qquad
 \sum_c\ell_{t,c}\le a'_t,\qquad
 \sum_c\ell_{t,c}\equiv a'_t\pmod{P_j(x)}.
\]
Write $n_j(x)=(a'_t-\sum_c\ell_{t,c})/P_j(x)$.
The terminal parameters satisfy $0\le e_{j,i}<rQ_j$ and
$e_{j,i}\equiv n_j(f_{j,i})\pmod{Q_j}$.
For allowed segment monomials $m_j$, the final exponent at a frame plane coordinate has the form
\[
 N_{t,c}=\ell_{t,c}+P_j(x)\bigl(m_j(x,c)+Q_jm_{j+1}(x,c)\bigr),
 \qquad m_j(x,c)<Q_j. \tag{LA.3}
\]
Indices and letters are translated periodically when necessary.
Reduction modulo $P_j(x)$ recovers $\ell$; division and reduction modulo $Q_j$ recover the current exponent and then the old exponent.
At a nonframe position, or a normal coordinate $c\ge r$, the old term is zero, so division alone suffices.
Outside $TT$ all exponents vanish.
Finally the current-frame row sums recover $e_{j,i}$.
Thus the exponent map is injective across \emph{all} label and terminal-parameter families, proving \texttt{RestLamInj}.

\paragraph{LA5. Exact row degrees, including natural subtraction.}
The prescribed degrees $b_j=b_{\ell,e,j}$ are $n_j(x)$ at intermediate letters, $e_{j,i}$ at current-frame letters, and
$(n_{j-1}(x)-e_{j-1,i})/Q_{j-1}$ at old-frame letters.
The last numerator is nonnegative: RS3 gives
\[
 \frac{a'_t}{P_j(x)Q_j}\ge\frac{36}{5}Yk^2N_s\ge k+r+1.
\]
Since $\sum_c\ell_{t,c}<kP_j(x)$, it follows that
$n_j(x)>Q_j(k+r+1)-k\ge rQ_j>e_{j,i}$.
The defining congruence therefore gives the exact identity
\[
 e_{j,i}+Q_j\frac{n_j(x)-e_{j,i}}{Q_j}=n_j(x).
\]
Summing (LA.3) over coordinates now gives $a'_t$ at every frame position; the intermediate positions give the same result directly.
This proves \texttt{RestLamDeg}, without importing the source's full \texttt{RestPad} argument.

\paragraph{LA6. The resulting dimension inequality.}
Choose the residue-independent bases from LA3 for every $(\ell,e,j)$.
LA4 sends distinct tensor-product monomials to distinct final monomials, so the corresponding product map is injective on the direct sum over $(\ell,e)$.
LA5 puts every product in the actual coarse generator set.
Consequently
\[
 \boxed{\displaystyle
 \sum_{\ell\in LA}\sum_{e\in EE(\ell)}
 \prod_{j=0}^{J-1}\operatorname{rank}_O M_j(b_{\ell,e,j})
 \le\dim_{\mathbb f}\overline\Lambda.} \tag{LA.4}
\]
This is the complete \texttt{RestLam} assembly. Its proof does not estimate the summands; their abundance belongs to the separate segment and residue-family counts.
